\chapter{长上下文——让模型阅读整本书}
\label{ch:long-context}
%% ============================================================

早期的 Transformer 模型只能处理 512 或 2048 个 token——
大约一两页文字。这严重限制了模型的实用性：
它无法阅读一篇长论文，无法理解一个完整的代码仓库，
更无法在长篇对话中保持一致性。

2025--2026 年，长上下文能力取得了巨大突破。
下表展示了截至 2026 年初主要模型的上下文窗口对比：

\begin{center}
\begin{tabular}{lcll}
\hline
\textbf{模型} & \textbf{上下文窗口} & \textbf{位置编码} & \textbf{发布时间} \\
\hline
LLaMA 4 Scout     & 10M tokens  & iRoPE       & 2025.04 \\
Gemini 2.5 Pro    & 1M tokens   & 未公开       & 2025.03 \\
Gemini 3 Pro/Flash & 1--2M tokens & 未公开      & 2025--2026 \\
Claude Sonnet 4.5 / Opus 4.5 & 200K--1M tokens & 未公开 & 2025--2026 \\
GPT-5 系列        & 400K tokens & 未公开       & 2025--2026 \\
DeepSeek-V3.2     & 1M tokens   & YaRN + MLA   & 2026.02 \\
Kimi K2.5          & 256K tokens & RoPE 变体    & 2026.01 \\
Qwen3-Next        & 262K--1M tokens & YaRN 变体 & 2026 \\
GLM-4.7           & 200K tokens & RoPE 变体    & 2025.12 \\
Jamba 1.5-Large   & 256K tokens & RoPE + SSM   & 2024.08 \\
DeepSeek-V3       & 128K tokens & YaRN + MLA   & 2024.12 \\
LLaMA 3.1 405B    & 128K tokens & RoPE + 频率调整 & 2024.07 \\
Qwen 2.5          & 128K--1M tokens & YaRN    & 2024--2025 \\
\hline
\end{tabular}
\end{center}

LLaMA 4 Scout 以 10M token（约 15,000 页）
创下了公开模型的最长上下文记录。
Google 的 Gemini 系列稳定支持 1--2M token。
Anthropic 的 Claude 系列在 2025--2026 年
将上下文扩展到了 200K--1M token。
OpenAI 的 GPT-5 系列达到了 400K token。

值得注意的是，\textbf{1M token 上下文在 2026 年初已成为事实上的标准}。
DeepSeek 在 2026 年 2 月 11 日悄然将 V3.2 的上下文窗口
从 128K 升级到 1M——没有发布公告，
用户在使用中自行发现了这一变化。
Kimi K2.5 支持 256K 原生上下文，
Qwen3-Next 原生支持 262K 并可扩展至 1M，
GLM-4.7 支持 200K 上下文和 128K 最大输出长度。
这意味着长上下文竞争的焦点
已经从``能否支持 1M''转向``能否有效利用 1M''。

但这些数字背后隐藏着艰巨的工程和科学挑战——
让模型「看到」长文本和让模型「理解」长文本，
完全是两回事。
正如 Chroma 的 context rot 研究~\cite{chroma2025contextrot} 所揭示的：
\textbf{标称上下文长度和有效上下文利用率之间存在巨大鸿沟。}

\section{为什么长上下文很难？}

长上下文的困难不是单一的，而是四重挑战的叠加。

\subsection{计算复杂度：$O(n^2)$ 的诅咒}

注意力机制的计算复杂度是 $O(n^2)$——
序列长度翻倍，计算量翻四倍。
让我们把这个抽象的复杂度换算成具体的数字：

处理 4K token 时，每一层注意力需要计算
$4{,}096 \times 4{,}096 = 16.8M$ 个注意力分数对。
处理 128K token 时，这个数字变成
$131{,}072 \times 131{,}072 \approx 16.4B$——
\textbf{增长了近 1000 倍}。
如果模型有 80 层（如 LLaMA 70B），
那就是 $80 \times 16.4B \approx 1.3$ 万亿次注意力分数计算。

即使有 FlashAttention~\cite{dao2022flashattention} 这样的优化
（将注意力的 IO 复杂度从 $O(n^2)$ 降到 $O(n^2/M)$，
其中 $M$ 是 SRAM 大小），
计算量本身的 $O(n^2)$ 增长是无法改变的。
128K token 的计算量就是 4K token 的 $1024\times$，
这是数学决定的，没有捷径。

\subsection{内存瓶颈：KV Cache 的爆炸}

比计算更致命的是内存。
在推理时，Transformer 需要缓存所有已处理 token 的
Key 和 Value 向量（称为 KV cache），
以便后续 token 做注意力计算。

KV cache 的大小可以这样估算：
\begin{equation}
  \text{KV cache} = 2 \times n_{\text{layers}} \times n_{\text{kv\_heads}}
  \times d_{\text{head}} \times n_{\text{tokens}} \times \text{bytes per element}
  \label{eq:kv-cache-size}
\end{equation}
其中因子 2 来自 Key 和 Value 两组向量，
$n_{\text{kv\_heads}}$ 是 KV 头的数量（注意不是 query 头的数量——
在使用 GQA 的模型中两者不同）。

LLaMA 70B 使用 GQA（Grouped Query Attention），
有 64 个 \textbf{query 头}但只有 8 个 \textbf{KV 头}
（每 8 个 query 头共享一组 KV 投影）。
因此 KV cache 计算直接使用 $n_{\text{kv\_heads}} = 8$。

对于 LLaMA 70B（80 层，8 个 KV heads，head 维度 128）
处理 128K token，使用 BF16（2 bytes）：
\[
  2 \times 80 \times 8 \times 128 \times 131{,}072 \times 2
  \approx 21 \text{ GB}
\]
仅一个用户的一次 128K 请求，
光是 KV cache 就要占用约 21\,GB——
超过单块 H100（80\,GB）显存的四分之一，
模型权重还没算。
如果不使用 GQA（即 64 个 KV heads），
KV cache 将膨胀到约 171\,GB，远超单块 GPU 容量。

\begin{warnbox}[KV Cache 是长上下文推理的真正瓶颈]
很多人以为长上下文的瓶颈是计算。
实际上，对于推理来说，\textbf{内存才是硬约束}。
计算可以通过更多 GPU 和更长时间来解决，
但 KV cache 必须驻留在 GPU 显存中，
而单块 GPU 的显存是有限的。
这就是为什么 KV cache 的压缩和分布式管理
是长上下文推理工程中最核心的问题。
\end{warnbox}

\subsection{训练--推理的长度失配}

第三个挑战更加根本：
模型在训练时只见过一定长度的序列，
当推理时序列长度超过训练长度时会发生什么？

RoPE~\cite{su2021rope} 的位置编码依赖于旋转角度
$\theta_i = 10000^{-2i/d}$。
在训练长度内（比如 4K），
模型见过的最大旋转角度是 $4096 \times \theta_i$。
当推理时序列长度变成 32K，
旋转角度变为 $32768 \times \theta_i$——
这些角度值是模型从未在训练中遇到过的。

直觉上，这就像一个只学过 0--100 范围内加法的学生
被要求做 0--1000 范围的加法。
虽然底层规则相同，但具体的数值模式是「未见过」的，
注意力分数变得不可靠——
这就是所谓的「外推失败」（extrapolation failure）。

低频维度受影响最大：
它们的旋转角度在训练范围内变化缓慢（不到一整圈），
但在超长序列中可能转过好几圈，
进入完全陌生的角度空间。
高频维度则相对安全，
因为它们在训练范围内已经转过了很多圈，
多转几圈不会带来本质变化。

\subsection{Lost in the Middle：中间信息的遗忘}

第四个挑战来自模型的注意力分配模式。
2023 年 Liu 等人的研究揭示了一个令人不安的现象：
当输入序列很长时，
模型对序列\textbf{开头}和\textbf{结尾}的信息关注度最高，
而对\textbf{中间}部分的信息检索准确度显著下降——
下降幅度超过 30\%。

% NEW REF: liu2023lost = Liu et al., "Lost in the Middle: How Language Models Use Long Contexts", arXiv:2307.03172, 2023

这种「U 型曲线」注意力模式有深层原因。
在自回归训练中，
每个 token 的损失只取决于它之前的 context。
序列开头的 token 是所有后续 token 的 context，
因此它们的信息被反复强化。
序列结尾的 token 距离当前预测位置最近，
在注意力分数中自然占优势（近因效应）。
中间的 token 则两头不沾——
既不是「总是被参考」的锚点，
也不是「最近看到」的新鲜信息。

这个发现意味着：仅仅增加上下文窗口是不够的。
如果关键信息恰好落在长序列的中间部分，
模型可能会「看到但不记住」——
就像一个学生虽然读了整本教科书，
但只记得开头和结尾的内容。

\section{位置编码的扩展方法}

面对训练--推理的长度失配问题，
研究者们提出了一系列巧妙的位置编码扩展方法。
这些方法的核心思想都是：
让超长序列中的位置编码落入模型训练时见过的值域范围内。

\subsection{线性位置插值}

最直接的方法是\textbf{线性插值}：
如果要将 4K 上下文扩展到 32K（$8\times$），
就把所有位置索引除以 8。
原来位置 32,768 的 token 变成位置 4,096——
刚好是训练时见过的最大位置。

% NEW REF: chen2023extending = Chen et al., "Extending Context Window of Large Language Models via Positional Interpolation", arXiv:2306.15595, 2023

数学上，将位置 $m$ 替换为 $m/s$，
其中 $s = L_{\text{target}} / L_{\text{train}}$ 是扩展比例。
RoPE 的旋转角度变为 $(m/s) \cdot \theta_i$，
确保所有角度都落在训练时见过的范围内。

为什么这个简单的技巧有效？
因为它保持了\textbf{相对位置关系}不变——
如果 token A 和 token B 的相对位置差是 $\Delta$，
插值后变成 $\Delta / s$，但两者的差值比例没变。
模型学到的「距离越远注意力越弱」的模式仍然适用。

但线性插值有明显的局限：
当扩展比例 $s$ 很大时（比如 $32\times$ 以上），
相邻 token 的位置差变得极小
（从 1.0 变成 $1/32 = 0.03125$），
模型难以区分紧邻的 token 的位置关系——
局部分辨率严重下降。
这就像把一张高清照片压缩到很小：
全局结构保留了，但细节模糊了。

\subsection{NTK-aware 插值}

NTK-aware 插值的核心洞察是：
\textbf{不应该对所有频率维度做相同比例的缩放}。

回忆 RoPE 的频率：$\theta_i = b^{-2i/d}$（$b = 10000$）。
高频维度（$i$ 小）负责编码\textbf{局部}位置关系
（相邻 token 之间的差异），
低频维度（$i$ 大）负责编码\textbf{全局}位置关系
（文档开头 vs 文档结尾）。

线性插值对所有频率一视同仁地除以 $s$，
结果高频维度——负责局部关系的维度——被不必要地压缩了。
NTK-aware 插值只调大 base 值 $b$：
将 $b = 10000$ 替换为 $b' = 10000 \cdot s^{d/(d-2)}$。

这产生了一种非均匀缩放效果：
\begin{itemize}[noitemsep]
  \item \textbf{高频维度}几乎不受影响——
        它们的频率本来就高，稍微调整 base 对它们影响很小。
        局部位置分辨率得以保留。
  \item \textbf{低频维度}被显著缩放——
        让它们能编码更远的位置关系。
\end{itemize}

一个有趣的类比：这就像数字系统中的进制变化。
在十进制中，三位数能表示 0--999。
如果我们把进制从 10 变成 100，
三位数就能表示 0--$100^3 = 999{,}999$——
扩大了 1000 倍的范围，
但低位的分辨率几乎不变。

NTK-aware 插值的最大优势是\textbf{不需要微调}——
在中等扩展比例（$4\times$ 左右）下，
直接修改 base 值就能在推理时扩展上下文，
质量损失很小。但在高扩展比例下仍会退化。

\subsection{YaRN：优雅的频率操控}

YaRN（Yet Another RoPE extensioN）是当前最流行的上下文扩展方法。
它的核心洞察是：RoPE 的不同频率维度应该被\textbf{区别对待}，
而且不仅要调整频率，还要调整注意力的温度。

% NEW REF: peng2023yarn = Peng et al., "YaRN: Efficient Context Window Extension of Large Language Models", arXiv:2309.00071, 2023

RoPE 的频率 $\theta_i = 10000^{-2i/d}$ 从高频（$i$ 小）到低频（$i$ 大）。
YaRN 将这些频率分为三组，并定义了一个关键概念——
频率维度的「波长」$\lambda_i = 2\pi / \theta_i$。
如果 $\lambda_i$ 远大于训练长度 $L$，
说明该维度在训练中连一整圈都没转完，外推风险最高。
如果 $\lambda_i$ 远小于 $L$，
说明该维度已经转过很多圈，外推风险最低。

基于这个分析，YaRN 将频率维度分为三组：
\begin{itemize}[noitemsep]
  \item \textbf{高频维度}（$\lambda_i < L$，对应短距离模式）：
        不做任何修改——这些维度已经能处理局部关系，不需要调整。
        它们在训练中转过了很多圈，
        外推到更长序列时只是多转几圈而已。
  \item \textbf{低频维度}（$\lambda_i > s \cdot L$，对应长距离模式）：
        按上下文扩展比例进行激进的频率缩放——
        比如要将 4K 扩展到 128K（$32\times$），
        这些维度的频率除以 32。
        它们必须被压缩到训练范围内，否则外推完全不可靠。
  \item \textbf{中频维度}（$L \leq \lambda_i \leq s \cdot L$）：
        进行渐进式的 NTK-aware 插值——
        在不缩放和完全缩放之间做平滑过渡。
        过渡函数是一个 ramp（线性插值因子），
        避免了频率空间中的突变。
\end{itemize}

此外，YaRN 引入了\textbf{注意力温度缩放}因子 $\sqrt{t}$。
长序列中注意力分数的方差会增大
（因为更多的 token 参与 softmax 竞争），
导致注意力分布变得过于「尖锐」——
集中在少数 token 上，忽略其余信息。
温度因子 $\sqrt{t}$ 相当于在 softmax 之前
将注意力 logits 除以 $\sqrt{t}$（$t > 1$），
降低 softmax 的「锐度」，
让注意力分布更加平滑，
使模型能更均匀地关注长序列中的信息。

通过这些手段，YaRN 只需在原始预训练数据的极小比例上做微调
（通常 $\sim$0.1\% 的训练数据，训练 $\sim$400 步），
就能将上下文窗口扩展 $32\times$ 甚至更多。
Qwen2.5、DeepSeek 等模型均使用 YaRN 变体进行上下文扩展。

\subsection{Dynamic NTK：推理时自适应}

Dynamic NTK 是一种在推理时根据实际序列长度
\textbf{动态调整}插值比例的方法。

核心思想很简单：不预设固定的扩展比例，
而是在推理时检查当前序列长度 $n$——
如果 $n \leq L_{\text{train}}$，不做任何修改；
如果 $n > L_{\text{train}}$，
将插值比例设为 $s = n / L_{\text{train}}$，
然后应用 NTK-aware 缩放。

这种方法的优势是\textbf{完全免微调}——
不需要任何额外训练，只需修改推理代码。
对于中等程度的扩展（$2\times$--$4\times$），
Dynamic NTK 的质量损失可以接受。
但对于大比例扩展（$8\times$ 以上），
质量下降明显，因为模型没有在长序列上做过任何适应。

Dynamic NTK 在实际应用中常被用作一种「应急方案」——
当用户的输入偶尔超过训练长度时，
它能优雅地降级，而不是直接崩溃。

\section{iRoPE：Meta 的交替架构}

LLaMA 4 引入了更激进的方案——iRoPE（interleaved RoPE）。
它不是在位置编码的数值上做文章，
而是从\textbf{架构层面}重新思考位置编码的角色。

\subsection{核心设计：交替堆叠}

iRoPE 交替堆叠两种注意力层：
\begin{itemize}[noitemsep]
  \item \textbf{RoPE 层}：使用旋转位置编码，
        捕捉局部上下文和精确的相对位置关系。
        这些层知道「哪个 token 在前，哪个在后」。
  \item \textbf{NoPE 层}（No Positional Encoding）：
        完全没有位置编码，注意力只基于语义内容。
        这些层天然不受序列长度限制——
        因为它们根本不编码位置信息。
\end{itemize}

为什么这种交替架构有效？
关键洞察是：\textbf{不是所有注意力模式都需要位置信息}。

考虑两种典型的注意力模式：
一种是「这个代词指代的是三句话前的名词」——
这需要精确的位置感知（RoPE 层的工作）。
另一种是「这段话在讨论量子物理的概念」——
这只需要语义匹配，不关心具体位置（NoPE 层的工作）。

NoPE 层的存在使得模型对长序列有天然的适应性——
它们关注的是「这个 token 的内容是什么」
而非「这个 token 在哪个位置」。
当序列从 8K 扩展到 1M 时，
NoPE 层完全不受影响，
只有 RoPE 层需要做位置编码的扩展。
这等于把需要外推的层数减少了一半。

\subsection{LLaMA 4 的 iRoPE 实现细节}

LLaMA 4 系列包含三个模型，
全部采用 iRoPE + MoE 架构：

\begin{center}
\begin{tabular}{lccccc}
\hline
\textbf{模型} & \textbf{总参数} & \textbf{活跃参数} & \textbf{专家数} & \textbf{上下文} \\
\hline
Scout   & 109B  & 17B  & 16    & 10M tokens \\
Maverick & 400B & 17B  & 128   & 1M tokens \\
Behemoth & 2T   & 288B & 16    & --- (训练中) \\
\hline
\end{tabular}
\end{center}

Scout 的 10M token 上下文窗口（约 15,000 页文本）
是截至 2025 年 4 月所有公开模型中最长的。
这个惊人的长度正是得益于 iRoPE 的架构设计：
NoPE 层提供了「免费」的上下文扩展空间，
RoPE 层通过频率调整处理位置敏感的注意力模式。

在 LLaMA 4 的具体实现中~\cite{meta2025llama4}，
iRoPE 的交替模式并非简单的 1:1 交替。
Meta 发现不同的交替比例适合不同的模型规模——
Scout 和 Maverick 使用了针对各自专家数量和模型规模
优化的交替策略。

iRoPE 还与 LLaMA 4 的其他架构创新协同工作：
\begin{itemize}[noitemsep]
  \item \textbf{早期融合}（Early Fusion）：
        LLaMA 4 是原生多模态模型，
        文本、图像和视频 token 在模型早期就被融合处理。
        NoPE 层对跨模态注意力特别有价值——
        「这张图片展示的概念」和「文本描述的概念」
        之间的匹配是纯语义的，不需要位置信息
  \item \textbf{MoE 稀疏激活}：
        每个 token 只激活少量专家，
        使得 Scout 的 109B 总参数中仅有 17B 活跃——
        这让 10M 上下文在 8$\times$H100 上可运行
\end{itemize}

\begin{keybox}[iRoPE 的训练流程]
LLaMA 4 的长上下文训练分三个阶段：
\begin{enumerate}[noitemsep]
  \item \textbf{预训练}：在 8K 上下文长度上完成主体训练。
        这个阶段消耗了绝大部分计算。
        Behemoth 作为「教师模型」在这个阶段训练，
        Scout 和 Maverick 通过蒸馏获得高质量初始化。
  \item \textbf{中等长度扩展}：调整 RoPE 层的频率参数，
        将上下文扩展到 128K--256K，
        在长文档数据上做少量继续训练。
  \item \textbf{超长扩展}：进一步扩展到 1M--10M。
        NoPE 层在这个阶段几乎不需要适应，
        RoPE 层通过 YaRN 风格的频率操控进行扩展。
\end{enumerate}
LLaMA 4 Scout 通过 iRoPE 达到了 10M token 的上下文窗口——
这在纯 RoPE 架构中几乎不可能实现。
\end{keybox}

\subsection{iRoPE 的推理优化}

iRoPE 的另一个优势是推理效率。
NoPE 层可以使用更激进的 KV cache 压缩策略
（比如丢弃远距离的 KV），
因为它们的注意力模式基于语义相似度而非位置，
远距离的 token 只要语义不相关就不重要。
而 RoPE 层则保留完整的位置信息，
确保精确的位置感知不受影响。

vLLM 在 v0.8.3 版本中已原生支持 LLaMA 4 的 iRoPE 架构。
实测显示，使用 8$\times$H100，
vLLM 可以为 Scout 提供 1M 上下文的推理服务，
为 Maverick 提供约 430K 上下文的推理服务。

\begin{warnbox}[iRoPE 的局限与争议]
LLaMA 4 在发布后收到了一些社区反馈：
虽然 10M token 的上下文窗口在 NIAH 等检索测试上表现优秀，
但在需要\textbf{深度推理}的长上下文任务上，
实际有效利用的上下文长度远低于标称值。
这再次验证了「标称上下文长度 $\neq$ 有效上下文长度」的原则。
此外，Behemoth 模型截至 2026 年初仍未正式发布，
社区对其实际性能持审慎态度。
\end{warnbox}

\section{Ring Attention：跨 GPU 处理超长序列}

即使解决了位置编码的问题，
超长序列的 KV cache 仍然可能超过单块 GPU 的内存容量。
Ring Attention 的解决方案是
将长序列\textbf{分割到多块 GPU 上}。

\subsection{核心思想：环形传递}

Ring Attention 的灵感来自一个简单的观察：
标准注意力要求每个 query 都能访问所有的 key 和 value，
但没人说所有 KV 必须\textbf{同时}在同一块 GPU 上。

设备排列成环形拓扑。
每块 GPU 持有序列的一个 chunk（包含该 chunk 的 Q、K、V），
然后执行以下步骤：
\begin{enumerate}[noitemsep]
  \item 每块 GPU 计算自己持有的 Q 与\textbf{当前本地} KV 的注意力。
  \item \textbf{同时}，每块 GPU 将自己的 KV block
        发送给环上的下一块 GPU，
        并从前一块 GPU 接收新的 KV block。
  \item 收到新的 KV block 后，计算 Q 与这些新 KV 的注意力，
        将结果与之前的注意力输出做\textbf{在线聚合}
        （online softmax aggregation——
        不需要存储完整的注意力矩阵，
        只需要维护 running sum 和 running max）。
  \item 重复步骤 2--3，直到 KV block 在环上转了一整圈，
        每块 GPU 都看到了所有的 KV。
\end{enumerate}

关键的工程技巧在于步骤 2：
\textbf{通信与计算重叠}。
当 GPU 在计算当前 KV block 的注意力时，
下一个 KV block 正在通过网络传输过来。
只要计算时间大于传输时间——
在序列足够长的情况下通常成立——
传输延迟就被完全隐藏了。

\subsection{内存与扩展性分析}

Ring Attention 将每块 GPU 的内存需求
从 $O(n)$（存储完整 KV cache）降为 $O(n/P)$
（$P$ 是 GPU 数量），
理论上支持无限长的序列——
只要有足够多的 GPU。
PyTorch 已经将 Ring Attention 原生集成到
ContextParallel API 中。

但实际的扩展性受两个因素限制：
\begin{itemize}[noitemsep]
  \item \textbf{通信带宽}：每个 KV block 需要在环上传递 $P-1$ 次。
        当 GPU 数量很多时，KV 的总通信量为 $O(n \cdot (P-1))$，
        网络带宽成为瓶颈。
        在 NVLink（900 GB/s）互连的节点内这不是问题，
        但跨节点（通常 100--400 Gb/s InfiniBand）时延迟显著增加。
  \item \textbf{收益递减}：超过一定长度后，
        模型的有效信息利用率下降（context rot 效应），
        更多的 GPU 和更长的上下文不一定带来更好的结果。
\end{itemize}

\subsection{Zig-Zag Ring Attention：负载均衡}

Ring Attention 有一个变体叫 \textbf{Zig-Zag Ring Attention}，
它解决了原始方法在自回归模型中的负载不均衡问题。
在因果注意力中（每个 token 只能关注之前的 token），
序列开头的 GPU 计算量最小（只有少量的 KV 对），
序列末尾的 GPU 计算量最大（需要关注所有之前的 token）。

如果将序列按顺序均匀分配给 $P$ 块 GPU，
第 1 块 GPU 的计算量约为第 $P$ 块的 $1/P$——
当 $P = 8$ 时，最轻的 GPU 只有最重的 $1/8$ 的工作量，
7/8 的时间在空等。

Zig-Zag 通过锯齿形的序列分割策略平衡各 GPU 的负载：
第一块 GPU 拿第 1、4、5、8 段序列，
第二块拿第 2、3、6、7 段……
使得每块 GPU 处理的总 token 数近似相等。
直觉上，每块 GPU 同时拥有序列开头（计算量小）
和序列末尾（计算量大）的 chunk，
两者的计算量互相平衡。

\textbf{Striped Attention} 是另一种变体，
它将序列按 token 级别交错分配
（第 $i$ 个 token 分配给第 $i \mod P$ 块 GPU），
实现了理论上完美的负载均衡，
但通信模式更复杂。

\section{长上下文训练的实践}

支持长上下文不是简单地「把序列长度调大」就行。
需要在数据、训练策略和工程三个维度同时发力。

\subsection{多阶段扩展训练}

通常需要分阶段训练：

\textbf{阶段一}：在标准上下文长度（如 4K--8K）上完成主体预训练，
消耗大部分训练 token（如 14T 中的 13.5T）。
这个阶段的重点是学习语言和知识，上下文长度不是关注点。

\textbf{阶段二}：逐步扩展上下文长度——
先扩展到 32K，训练少量数据适应；
然后扩展到 128K，再训练一批数据。
Qwen2.5 的长上下文版本甚至使用了三阶段扩展（32K → 128K → 1M）。
每次扩展时需要调整 RoPE 的频率参数（通过 YaRN 或类似方法）。

LLaMA 3~\cite{dubey2024llama3} 的长上下文训练提供了一个典型的工程实例：
预训练在 8K 上下文上完成，
然后通过六个子阶段逐步扩展到 128K——
每个子阶段将上下文长度翻倍（8K → 16K → 32K → ... → 128K），
同时逐步增大 RoPE 的 base 频率
（从 $5 \times 10^5$ 逐步增大到 $8 \times 10^6$）。

\textbf{阶段三}：在长上下文数据上做 SFT，
确保模型能够在长文本中遵循指令、
回答关于文档特定部分的问题、
以及保持长距离的一致性。

这种分阶段策略的原因是：
长序列的训练效率远低于短序列——
注意力的 $O(n^2)$ 复杂度意味着
一条 128K 的样本比 32 条 4K 的样本\textbf{慢得多}，
但包含的有效 token 数是一样的。
在短序列上完成大部分学习，
然后在长序列上做少量适应，是效率最优的策略。

\subsection{长上下文训练数据}

长上下文训练面临一个独特的数据挑战：
\textbf{大多数高质量文本都不长}。
网页平均长度约 2K--4K token，
新闻文章约 500--2000 token。
真正超过 32K token 的高质量文档并不多。

长上下文训练数据的来源：
\begin{itemize}[noitemsep]
  \item \textbf{书籍}：单本书通常在 50K--500K token，
        是最天然的长上下文数据源。
        但版权问题限制了可用数量。
  \item \textbf{代码仓库}：一个完整的代码仓库
        （包含所有文件及其依赖关系）
        可以轻松超过 100K token。
        代码的优势是结构清晰、
        有明确的跨文件引用关系。
  \item \textbf{学术论文}：单篇论文 8K--30K token，
        可以将相关论文串联起来作为长文档。
  \item \textbf{法律和财务文档}：合同、招股书、法规
        动辄数十万 token，是天然的长文本。
  \item \textbf{合成长文档}：将多个相关短文档
        按逻辑顺序拼接成长文档。
        例如将关于同一主题的 Wikipedia 页面拼接起来，
        或者在一组文档上生成跨文档的问答对。
        这种方法需要小心处理拼接边界，
        避免模型学到「文档之间无关联」的错误模式。
\end{itemize}

\subsection{长上下文 SFT 数据}

阶段三的长上下文 SFT 需要专门设计的指令数据。
典型的任务类型包括：
\begin{itemize}[noitemsep]
  \item \textbf{长文档总结}：给出一篇 50K+ token 的文档，
        要求模型生成不同粒度的摘要（一句话/一段/多页）。
  \item \textbf{多文档问答}：给出多篇文档，
        要求模型综合多个来源回答问题。
        这测试的是跨文档信息整合能力。
  \item \textbf{仓库级代码理解}：给出整个代码仓库，
        要求模型解释架构、定位 bug、
        或在保持全局一致性的前提下修改代码。
  \item \textbf{长距离事实一致性}：
        在对话的早期提供一组事实，
        在很久之后（上万 token 的对话之后）
        询问与这些事实相关的问题。
\end{itemize}

\subsection{推理工程优化}

SGLang 在 2026 年初发布的 Chunked Pipeline Parallelism
是长上下文推理的重要工程突破。
它将长序列的 prefill（首次处理整个输入序列的阶段）
分成多个 chunk，在不同 GPU 上以流水线方式并行处理。
在 DeepSeek-V3.1 上的实测显示，
处理 1M token 的输入延迟从 55.5 秒降至 10.5 秒——
降低了 81\%，相当于 $3.31\times$ 的吞吐提升。
这意味着用户可以在约 10 秒内完成对一整本书的「阅读」，
极大地提升了长上下文模型的实用性。

\section{长上下文的实际应用场景}

长上下文能力开启了许多此前不可能的应用场景：

\textbf{整仓库代码理解}：
将一个完整的代码仓库（数万行代码）放入上下文，
让模型理解整体架构后再回答具体问题或做修改。
这比传统的 RAG（只检索相关片段）
能获得更全面的上下文理解。
Cursor 和 GitHub Copilot Workspace 等开发工具
已经在利用这种能力——
将整个项目的代码、配置文件、README
一起放入上下文，让模型做出全局一致的修改。

\textbf{长文档分析}：
将一份数百页的法律合同或学术论文完整放入上下文，
让模型进行总结、比较或提取特定信息。
这避免了将文档切分成小块后丢失跨段落关联的问题。
在金融领域，分析师可以将一整份招股书
（通常 200--500 页）放入上下文，
直接询问「这家公司的主要风险因素和竞争优势分别是什么？」

\textbf{多文档对比分析}：
将多篇研究论文同时放入上下文，
让模型比较它们的方法、结果和结论。
这在文献综述和系统性综述中特别有价值——
传统做法需要人类研究者反复翻阅多篇论文，
而长上下文模型可以「同时阅读」所有论文并做对比。

\textbf{多轮对话的长期记忆}：
在一次持续数小时的编程对话中，
模型可以记住早期讨论的设计决策，
在后期做修改时保持一致性。
对于短上下文模型，长对话必须依赖
RAG 或摘要来维持「记忆」，
但这些方法不可避免地会丢失信息。
足够长的上下文窗口可以直接保留完整的对话历史。

\textbf{少样本学习}：
在上下文中放入大量的示例和说明，
让模型在不做任何微调的情况下学会新任务。
更长的上下文意味着可以提供更多、更多样的示例。
一个极端的例子：在上下文中放入整本编程教材，
然后要求模型按照教材的风格和约定写代码——
这比用几个 few-shot 示例效果好得多。

\textbf{长内容生成}：
生成长篇小说章节、技术文档、剧本等需要
在数万 token 的范围内保持情节、角色、
技术概念一致性的内容。
短上下文模型在长篇生成中常出现前后矛盾——
因为它「忘记了」前面写过什么。

但正如下一节将讨论的，
更长的上下文不等于更好的利用。
在实际应用中，如何组织和筛选放入上下文的信息
（context engineering）与上下文长度本身同样重要——
甚至更重要。

\section{Context Rot 与 NIAH：长上下文的质量评估}

在追求更长的上下文窗口时，
我们需要区分两个不同的能力维度：
\textbf{检索}（retrieval）和\textbf{理解}（reasoning）。
2025 年的研究表明，这两个维度之间的差距
远比人们最初预期的更大。

\subsection{Needle in a Haystack 测试}

NIAH（Needle in a Haystack）是最广泛使用的长上下文评估方法。
测试方法很简单：在一段无关的长文本（haystack）中
插入一条关键信息（needle），
然后询问模型关于这条信息的问题。
通过改变 needle 的位置（开头/中间/结尾）
和 haystack 的长度，
可以绘制出模型在不同位置和长度下的检索准确率热力图。

NIAH 测试揭示了几个重要模式：
\begin{itemize}[noitemsep]
  \item 大多数现代模型在 NIAH 上表现优秀——
        即使在 128K 长度下，主流模型的检索准确率也接近 100\%。
  \item 但 NIAH 本质上只测试\textbf{检索}——
        模型能否找到特定信息。
        这是长上下文能力的\textbf{必要条件}，
        但远不是\textbf{充分条件}。
\end{itemize}

一个模型可以完美地通过 NIAH（找到 needle），
但在需要\textbf{综合}长上下文中多处信息做推理时表现糟糕。
比如「根据文档第 3 章和第 7 章的描述，
这两部分的结论是否矛盾？」——
这要求模型不仅找到两处信息，
还要理解并比较它们的含义。

\subsection{NIAH 的局限与更好的评估}

NIAH 的根本问题是它把长上下文简化成了一个搜索任务。
在真实场景中，用户很少需要「在长文档中找一句话」——
更常见的是需要\textbf{综合}、\textbf{推理}和\textbf{比较}。

更有挑战性的评估任务包括：
\begin{itemize}[noitemsep]
  \item \textbf{Multi-Needle NIAH}：
        在 haystack 中插入\textbf{多个}相关信息片段，
        要求模型综合所有片段才能回答问题。
        这测试的是信息\textbf{聚合}能力，
        而非单点检索
  \item \textbf{RULER}（2024）：
        包含多种长上下文推理任务——
        变量追踪、多跳推理、
        以及需要整合分散在上下文各处的信息的复合任务
  \item \textbf{长文档摘要质量}：
        给模型一篇 100K+ token 的文档，
        评估摘要是否遗漏了中间部分的关键信息。
        这直接测量 lost-in-the-middle 效应
  \item \textbf{BABILong}：
        基于 bAbI 任务改编的长上下文推理基准，
        需要在长文本中追踪实体状态和关系
\end{itemize}

\subsection{Context Rot：系统性退化}

2025 年 7 月，Chroma 发布了影响深远的
\textbf{context rot} 研究报告~\cite{chroma2025contextrot}，
对 18 个 frontier 模型（包括 GPT-4.1、Claude 4、
Gemini 2.5、Qwen3 等当时最新的模型）
进行了系统性的长上下文评估。
核心发现令人警醒：
\textbf{所有模型都随着输入长度的增加而性能下降，
即使远未达到标称的上下文窗口上限}。

具体发现：
\begin{itemize}[noitemsep]
  \item \textbf{U 型曲线持续存在}：
        模型对序列开头和结尾的信息关注度最高，
        中间部分的信息检索准确度下降超过 30\%。
        这被称为「lost in the middle」效应~\cite{liu2023lost}。
        尽管 2024--2025 年的训练改进缓解了部分问题，
        但 U 型模式并未消失——只是变得不那么陡峭
  \item \textbf{反直觉发现}：有逻辑结构的上下文
        比随机打乱的上下文\textbf{更难}检索——
        模型在有意义的干扰信息中更容易「迷路」。
        这意味着真实应用场景
        （放入完整文档而非随机文本）的难度
        比 NIAH 测试暗示的更高
  \item \textbf{推理 vs 检索的鸿沟}：
        在需要跨上下文推理的任务中，
        性能下降比纯检索任务更为严重。
        模型可能能找到两个相关段落，
        但无法正确地将它们联系起来做推理。
        对于简单的重复词计数任务，
        一些模型在输入增长到几千 token 时
        就开始出现可靠性问题
  \item \textbf{标称 vs 有效上下文}：
        实际测试显示，
        许多模型的「有效上下文利用率」（实际可靠使用的比例）
        仅为标称上下文窗口的 50--80\%。
        例如一个标称 200K 的模型，
        在约 100--160K token 后可靠性就开始下降
  \item \textbf{核心结论}：
        \textbf{Context engineering（如何组织上下文）
        比 context length（能塞多少）更重要}
\end{itemize}

\subsection{Context Engineering：从「更长」到「更智能」}

Anthropic 在 2025 年下半年
明确将 ``context engineering'' 定义为构建高效 AI Agent 的核心技能，
标志着行业思维从「追求更长的上下文窗口」
转向「更智能地利用有限的上下文预算」~\cite{anthropic2025contexteng}。

\begin{whybox}[Context Engineering 的核心原则]
基于 context rot 的研究和生产实践，
以下原则已成为行业共识：
\begin{itemize}[noitemsep]
  \item \textbf{重要信息前置}：
        把最关键的指令和信息放在上下文的开头，
        利用模型对开头信息的高关注度。
        系统提示词（system prompt）和核心约束应在最前面
  \item \textbf{冗余而非堆积}：
        对关键信息做适度冗余
       （比如在不同位置重复强调重要约束），
        而不是试图塞入尽可能多的不同信息。
        每个不必要的 token 都在\textbf{主动}降低其他 token 的注意力预算
  \item \textbf{结构化组织}：
        使用清晰的标题、分隔符和层次结构
        帮助模型定位和组织信息。
        XML 标签、Markdown 标题、编号列表
        都能显著提高信息检索的可靠性
  \item \textbf{即时上下文（Just-in-Time Context）}：
        不要在一开始就加载所有信息，
        而是在需要时才检索和注入——
        这对 AI Agent 尤为重要。
        一个执行 20 步任务的 Agent，
        在每一步只需要当前相关的上下文，
        而非全部 20 步的信息
  \item \textbf{上下文压缩（Compaction）}：
        对长时间运行的 Agent，
        定期将已处理的上下文压缩为简洁的摘要，
        丢弃不再需要的细节。
        这类似于人类的「工作记忆管理」
  \item \textbf{RAG + 长上下文的协同}：
        RAG 负责从海量文档中筛选相关内容，
        长上下文负责一次性理解所有筛选出的内容。
        两者互补而非替代。
        经验表明 RAG 检索 top-20 到 top-50 相关片段
        并放入 32K--128K 的上下文中，
        效果通常优于直接将整个文档塞入 1M 上下文
\end{itemize}
\end{whybox}

\section{稀疏注意力的新进展}
\label{sec:sparse-attention-2026}

2026 年初，多项研究在稀疏注意力方向取得了突破，
为长上下文推理提供了更高效的计算范式。
这些工作的共同目标是：
在不显著损失质量的前提下，
减少长序列中实际参与注意力计算的 token 数量。

\subsection{Token Sparse Attention}

Token Sparse Attention~\cite{tokensparse2026}（arXiv:2602.03216）
提出了一种\textbf{逐头 Q/K/V 压缩}策略。
不同于 NSA 在 token 级别做选择，
Token Sparse Attention 在每个注意力头内部
对 Q、K、V 向量进行独立的压缩，
根据每个头学到的重要性分数
动态决定哪些 token 的表示需要保留完整精度，
哪些可以被压缩或跳过。
这种逐头的精细化控制
避免了``一刀切''式的全局稀疏策略
在某些头上过度丢弃信息的问题。

\subsection{Double-P 分层稀疏注意力}

Double-P~\cite{doublep2026}（arXiv:2602.05191）
引入了\textbf{分层 top-$p$ 稀疏注意力}机制。
其核心思想是将稀疏选择分为两个层次：
第一层 top-$p$ 在粗粒度（token 块级别）上筛选相关区域，
第二层 top-$p$ 在精细粒度（单 token 级别）上
从候选区域中选择最相关的 token。
这种两级筛选策略
既保持了全局上下文的感知能力，
又避免了在完整序列上做精细排序的计算开销。

\subsection{DySCO：动态注意力缩放}

DySCO（Dynamic Scaling for Context Optimization）
~\cite{dysco2026}（arXiv:2602.22175）
提出了一种基于\textbf{检索头}（retrieval heads）
的动态注意力缩放机制。
研究发现，Transformer 中存在少数``检索头''——
这些头在长上下文任务中负责精确检索远距离信息。
DySCO 识别这些检索头后，
动态调整不同头的注意力计算范围：
检索头保持对完整上下文的注意力，
非检索头使用更激进的稀疏策略。
这种自适应策略在不同任务上
自动找到效率和质量的平衡点。

\subsection{``Lost in the Middle'' 的精确理论}

2026 年 3 月，``Lost in the Middle''
效应首次获得了严格的理论解释~\cite{lostmiddle2026}（arXiv:2603.10123）。
研究者从 Transformer 的因果注意力结构出发，
数学地证明了 U 型注意力分布是
自回归训练目标和 softmax 归一化
在有限精度下的\textbf{必然结果}——
而非模型容量不足或训练数据不足导致的偶然现象。
这一理论结果意味着：
在不改变基本架构的前提下，
仅靠增加训练数据或模型参数
\textbf{无法彻底消除} lost-in-the-middle 效应。
解决方案必须来自架构层面的创新——
如稀疏注意力、分层处理、
或下文将讨论的长期记忆机制。

\subsection{UltraLong-8B：从 128K 到 4M 的高效训练}

UltraLong-8B~\cite{ultralong2025}（arXiv:2504.06214）
证明了一个令人振奋的可能性：
通过精心设计的多阶段训练策略，
一个仅有 8B 参数的模型可以从 128K
高效地扩展到 \textbf{4M token} 的上下文长度。
关键技术包括渐进式 RoPE 频率调整、
分阶段长度扩展训练、以及
针对超长序列的数据采样策略。
这表明超长上下文并不必然需要超大参数模型——
训练方法论的创新可以让中等规模模型
也具备百万级上下文处理能力。

\section{长期记忆：超越传统上下文窗口}
\label{sec:long-term-memory}

2026 年，长上下文领域迎来了一个根本性的范式转变——
从``如何让上下文窗口更长''
转向``如何让模型具有真正的长期记忆''。
传统的上下文窗口无论多长，
都受限于一个根本约束：
窗口内的所有 token 必须参与注意力计算，
计算和内存成本与窗口大小成正比（甚至平方关系）。

长期记忆（Long-Term Memory, LTM）
的核心思想是将``记忆''与``计算''分离——
让模型拥有一个远大于上下文窗口的
持久化知识存储，
在需要时以极低的成本检索。

\subsection{DeepSeek Engram：条件记忆}

2026 年 1 月 12 日，DeepSeek 发表了
一篇可能重塑 LLM 架构范式的论文——
``Conditional Memory via Scalable Lookup:
A New Axis of Sparsity for LLMs''
~\cite{engram2026}（arXiv:2601.07372），
同样由 CEO 梁文锋联合署名。

\subsubsection{核心洞察：Transformer 缺少查找原语}

Engram 论文的出发点是一个深刻的观察：
\textbf{Transformer 架构缺少原生的知识查找原语}。
当前模型回忆事实知识时，
必须通过多层注意力和 FFN 的计算来``模拟''检索——
这是一种极其低效的方式。
就好比一个图书馆员
每次找书都不查目录，
而是把所有书从头到尾翻一遍。

Engram 模块通过引入一种
\textbf{现代化的 N-gram 嵌入}实现了 $O(1)$ 的查找操作。
给定输入的局部上下文（几个相邻 token），
Engram 直接从一个巨大的嵌入表中
查找对应的知识表示——
不需要经过注意力计算。

\subsubsection{U 型缩放定律}

论文最引人注目的发现之一是
关于稀疏模型容量分配的\textbf{U 型缩放定律}：
在给定的总参数预算下，
分配给 Engram 记忆的最优比例约为 \textbf{75\%}——
也就是说，一个高效的稀疏模型
应该将四分之三的容量用于存储可查找的知识，
仅用四分之一的容量用于计算。

这个发现挑战了现有模型的设计惯例——
当前 LLM 几乎将所有参数都用于计算
（注意力和 FFN 层），
没有专门的知识存储空间。

\subsubsection{部署优势}

Engram 的另一个重要特性是
其存储可以放在\textbf{主机内存}（host memory）上
而非昂贵的 HBM 上。
即使是 100B 参数规模的 Engram 存储，
放在主机内存中也仅带来\textbf{不到 3\% 的吞吐量损失}——
因为 Engram 的查找操作是简单的表查询，
不需要 GPU 的高速计算能力。
这意味着在相同的 GPU 硬件上，
可以``免费''地为模型添加
数十倍于 GPU 显存容量的知识存储。

DeepSeek 已在 GitHub 上以 Apache 2.0 许可
开源了 Engram 实现（\texttt{deepseek-ai/Engram}，
截至 2026 年 3 月获得 3,840 个 star）。
Engram 被认为是
DeepSeek 实现 1M 上下文窗口的关键赋能技术之一。

\begin{keybox}[Engram 与传统方法的对比]
Engram 的条件记忆与已有的长上下文方法是互补的：
\begin{itemize}[noitemsep]
  \item \textbf{vs. 长上下文窗口}：
        上下文窗口适合处理当前会话中的信息，
        Engram 适合存储模型需要``永久记住''的知识
  \item \textbf{vs. RAG}：
        RAG 需要外部检索系统和向量数据库，
        Engram 将检索内置到模型架构中，
        延迟更低且无需外部依赖
  \item \textbf{vs. KV cache 压缩}：
        KV cache 压缩减少了计算中间结果的存储，
        Engram 增加了独立于计算的知识存储
\end{itemize}
\end{keybox}

\subsection{Google Titans：神经记忆架构}

2025 年 12 月，Google 发表了
\textbf{Titans 架构}~\cite{titans2025}，
引入了\textbf{Neural Memory 模块}——
一种可在运行时自我更新的长期记忆机制，
其设计灵感来源于生物大脑的记忆系统。

Titans 的核心架构将模型分为三个组件：
\begin{itemize}[noitemsep]
  \item \textbf{短期注意力核心}（Short-Term Attention Core）：
        标准的 Transformer 注意力层，
        负责处理当前窗口内的精确上下文交互
  \item \textbf{长期神经记忆}（Long-Term Neural Memory）：
        一个在推理时持续更新的神经网络模块，
        能够将跨越数百万 token 的信息
        压缩存储为持久化的记忆表示。
        不同于 RNN 的隐状态，
        神经记忆通过自监督学习目标
        在运行时动态更新其参数——
        真正实现了``边推理边学习''
  \item \textbf{持久记忆}（Persistent Memory）：
        存储不随输入变化的任务知识，
        类似于模型的``先验知识''
\end{itemize}

Titans 的突破性在于：
它能够处理\textbf{数百万 token}的输入，
同时保持跨越数周甚至数月对话的记忆连续性。
Google 将其称为``AI 失忆症的终结''——
传统模型在每次对话结束后
就``忘记''了所有交互内容，
而 Titans 的神经记忆可以在对话间持续积累知识。

从计算效率的角度看，
Titans 结合了 RNN 的处理速度
和 Transformer 的表示精度——
神经记忆模块的更新操作在时间复杂度上
类似于 RNN 的状态更新（$O(1)$ per token），
而短期注意力核心保证了
局部上下文的精确处理。

\subsection{Agent 记忆系统的爆发}

随着 LLM Agent（智能体）从研究原型走向生产部署，
Agent 的长期记忆管理在 2026 年初
成为了研究的焦点领域。
与模型层面的长期记忆不同，
Agent 记忆系统关注的是
\textbf{在多次交互和长时间跨度中
管理和利用经验知识}。

\subsubsection{代表性工作}

\textbf{AriadneMem}：
面向``终身 Agent''的记忆架构，
能够将分散在不同时间和不同任务中的
相关证据``穿针引线''般连接起来，
并通过主动的状态更新机制
保持记忆的一致性和时效性。

\textbf{PersistBench}：
首个针对 Agent 长期记忆的安全性基准，
揭示了两个重要风险：
（1）\textbf{跨域泄漏}——Agent 在处理任务 A 时
积累的敏感信息可能在任务 B 中被无意泄露；
（2）\textbf{记忆诱导的谄媚}（memory-induced sycophancy）——
Agent 可能过度迎合记忆中记录的用户偏好，
而非提供客观的最优建议。

\textbf{Adaptive Memory Structures}：
根据上下文自适应地选择记忆类型
（工作记忆、情景记忆、语义记忆），
而非使用单一的记忆存储。
这种设计模仿了人类认知心理学中的
多存储记忆模型。

\textbf{MemoryLLM}：
将持久化记忆 token 通过自注意力机制
融合到模型的前向传播中。
记忆 token 作为额外的 KV 对
参与每一层的注意力计算，
无需修改模型架构即可实现记忆增强。

2026 年 3 月发表的综述论文
``Memory for Autonomous LLM Agents''
~\cite{agentmemory2026}（arXiv:2603.07670）
对这一快速发展的领域进行了全面梳理，
将 Agent 记忆系统分类为
读取、写入、管理和组织四个维度，
并指出\textbf{记忆安全}将成为
Agent 大规模部署前必须解决的关键问题。

\begin{warnbox}[长期记忆的开放问题]
长期记忆虽然前景广阔，但仍面临多项挑战：
\begin{itemize}[noitemsep]
  \item \textbf{记忆可靠性}：如何确保长期记忆中的信息
        不会随时间``腐化''或产生虚假关联？
  \item \textbf{遗忘机制}：有选择地遗忘过时信息
        与``不可学习''（machine unlearning）的
        隐私合规需求如何协调？
  \item \textbf{评估困难}：如何系统地评估
        跨越数天、数周时间跨度的记忆能力？
        现有的基准测试远未解决这个问题
  \item \textbf{安全风险}：PersistBench 已经表明
        长期记忆可能引入新的安全攻击面。
        记忆中毒（memory poisoning）
        可能成为新型攻击向量
\end{itemize}
\end{warnbox}

\begin{corebox}[长上下文的未来方向]
截至 2026 年初，长上下文领域有几个值得关注的研究方向：
\begin{enumerate}[noitemsep]
  \item \textbf{混合架构}：iRoPE（LLaMA 4）和 Transformer-SSM 混合模型
        通过架构创新扩展上下文，
        而非仅依赖位置编码技巧。
        SSM 层不需要 KV cache，
        天然适合超长上下文
  \item \textbf{稀疏注意力训练}：
        DeepSeek 的 NSA 证明了稀疏注意力可以端到端训练，
        不再需要全注意力预训练后蒸馏的两阶段流程。
        Token Sparse Attention、Double-P、DySCO 等后续工作
        在 2026 年初进一步丰富了稀疏注意力的技术工具箱
  \item \textbf{KV cache 压缩}：
        MLA（DeepSeek）、GQA、动态 KV 驱逐等技术
        持续降低长上下文推理的内存开销
  \item \textbf{长期记忆}：
        Engram（DeepSeek）和 Titans（Google）
        开辟了``计算与记忆分离''的新范式，
        Agent 记忆系统的爆发
        进一步将记忆管理推向工程实践
  \item \textbf{无限上下文}：
        Ring Attention + 流水线并行使得
        理论上可以处理任意长度的序列，
        瓶颈从「能否处理」转变为「能否有效利用」
  \item \textbf{Context engineering 工具化}：
        自动上下文管理工具（如智能截断、
        动态检索增强、注意力预算分配）
        可能比扩展上下文窗口本身更有实际价值
\end{enumerate}
\end{corebox}


%% ============================================================
%% NEW REFERENCES (to be added to refs.bib):
%% meta2025llama4 = Meta. "The Llama 4 Herd: The Beginning of a New Era of Natively Multimodal AI Innovation." Meta AI Blog, April 2025.
%% chroma2025contextrot = Hong, Kelly et al. "Context Rot: How Increasing Input Tokens Impacts LLM Performance." Chroma Technical Report, July 2025.
%% liu2023lost = Liu, Nelson et al. "Lost in the Middle: How Language Models Use Long Contexts." arXiv:2307.03172, 2023.
%% anthropic2025contexteng = Anthropic. "Context Engineering for AI Agents." Anthropic Blog, 2025.
%%
%% --- ADDED 2026-03-12 ---
%% engram2026 = DeepSeek-AI. "Conditional Memory via Scalable Lookup: A New Axis of Sparsity for LLMs." arXiv:2601.07372, January 2026.
%% titans2025 = Google Research. "Titans: Learning to Memorize at Test Time." arXiv:2501.00663, December 2025.
%% tokensparse2026 = "Token Sparse Attention: Per-Head Q/K/V Compression for Long-Context Attention." arXiv:2602.03216, February 2026.
%% doublep2026 = "Double-P: Hierarchical Top-p Sparse Attention." arXiv:2602.05191, February 2026.
%% dysco2026 = "DySCO: Dynamic Attention-Scaling via Retrieval Heads." arXiv:2602.22175, February 2026.
%% lostmiddle2026 = "Lost in the Middle: An Exact Theory of Transformer Position Bias." arXiv:2603.10123, March 2026.
%% ultralong2025 = "UltraLong-8B: Efficient Training from 128K to 4M Tokens." arXiv:2504.06214, 2025.
%% agentmemory2026 = "Memory for Autonomous LLM Agents: A Comprehensive Survey." arXiv:2603.07670, March 2026.
%% ============================================================
